\chapter{Optimal Strategies as a Markov Decision Process}
\label{sec:mdp}

Rewriting the recursion as a MDP so it can take into account rescheduling options
including for things like temporal redundancy and reactive rescheduling.

\section{One-PE Formulation}
\label{sec:mdp:onepe}

\noindent{\textbf {State --}}
A state is a pair $(C, n)$, where $n$ is the number of time steps until
the deadline and $C$ is the set of all results reached so far. An element
$(i, x, y) \in C$ denotes an input of $x$ bytes and error level $y$ available
as input to task $T_i$.
Initially $C = \{(1, x_0, y_0)\}$ and $n = H$.  The output error so far is:
$$
y_{\min} (C) = \min\big(\efail, \min_{(M+1, x, y) \in C} y\big).
$$

\noindent{\textbf {Actions --}}
At time zero and after each success or failure, execute any task $T_i$ on
any of its available inputs with any execution mode $j$, given that it
can complete by the deadline and would produce a new output:
\begin{multline*}
A(C, n) = \{ (i, x, y, j) : (i, x, y) \in C, i \leq M, m^i_j \leq n,\\
(i+1, d^i_j(x,y), e^i_j(x,y)) \notin C \}.
\end{multline*}

\noindent{\bf Transitions --}
For action $(i, x, y, j)$:
\begin{compactitem}
\item \textbf{success, with probability $\psuc{m^i_j}$} \\
$(C \cup \{(i+1, d^i_j(x,y), e^i_j(x,y))\}, n - m^i_j)$
\item \textbf{failure during time step $u$, with probability $p_u$}\\ 
$(C, \max(0, n - u - 1 - \lceil R/\deltat \rceil))$.
\end{compactitem}

\noindent{\textbf {Optimal expected error --}}
Let $E(C,n)$ be the optimal expected error from state $(C,n)$.  If
$A(C,n) = \emptyset$ then $E(C,n) = y_{\min}(C)$. Otherwise:
\begin{multline*}
E(C, n) = \min_{(i,x,y,j) \in A(C,n)} \big[\\
\psuc{m^i_j} \times E(C \cup \{(i+1, d^i_j(x,y), e^i_j(x,y))\}, n - m^i_j) + \\
 \sum_{u=0}^{m^i_j - 1} p_u E(C, \max(0, n - u - 1 - \lceil R/\deltat \rceil)) \big].
\end{multline*}
The optimal expected error for the application is $E(\{(1, x_0, y_0)\}, H)$.

\section{Multi-PE Formulation}
\label{sec:mdp:multipe}

\noindent{\textbf {State --}}
A state is a triple $(C, n, \sigma)$. Set $C$ is shared by all PEs, and
$\sigma = (\sigma_1, \ldots, \sigma_N)$ gives the possible statuses of each PE:
\begin{compactitem}
\item available
\item executing task $T_i$ on input $(x,y)$ with mode $j$
\item restarting after a failure
\end{compactitem}